\subsection{The smallest basis rank with an obstruction}\label{sec:minimal-basis-rank}
The two-nonunit output relation in Forsberg's obstruction \citep[Corollary 13]{Forsberg} yields a failure of effective lifting already present for a free coefficient module of rank three. Here ``rank'' refers to the number of named basis vectors, including the unit; it is distinct from the rank-three nonstandard algebra, whose displayed basis has ten elements.

\paragraph{Proof of Theorem~\ref{thm:minimal-basis-rank-main}.}
\begin{proof}
First recall explicitly that $\pi$ is surjective on individual coefficients. At a present weight with integer trace $n>0$, choose $n$ copies of the trivial sign; for $n<0$, choose $-n$ copies of $\eps$; and for a present zero choose $1+\eps$. At an absent weight choose zero. A coefficient has finite support, so this constructs an element of $\mathcal E$ with the prescribed image in $\mathcal S$.

Basis rank one with strict unit has its unique unit multiplication table, which lifts. In basis rank two, write the basis as $1,x$. Its entire multiplication table is determined by
\[
 x^2=\alpha\,1+\beta x.
\]
Choose any coefficient lifts $\widetilde\alpha,\widetilde\beta\in\mathcal E$ and retain the strict unit. This table is associative. By bilinearity it suffices to check basis triples; those containing $1$ are immediate, and the only remaining triple satisfies
\[
 (x^2)x=\widetilde\alpha x+
 \widetilde\beta(\widetilde\alpha\,1+\widetilde\beta x)
 =x(x^2),
\]
using commutativity of the coefficient semiring. Hence every basis-rank-two table lifts.

For rank three, let $e$ denote the old supported zero at Laurent weight zero. On the free $\mathcal S$-semimodule with basis $1,a,b$, retain the strict unit and define
\[
 a^2=e a,\qquad ab=ba=e a+e b,\qquad b^2=e a+e b.
\]
All omitted coefficients are the unsupported zero. Since $e^2=e$ and $e+e=e$, the triples involving two $a$'s and one $b$, or one $a$ and two $b$'s, evaluate to $e a+e b$ in either parenthesization. The triples $a,a,a$ and $b,b,b$ are associative by commutativity, and strict-unit triples are immediate. Bilinearity now proves associativity on the whole free semimodule. Reflection sends every displayed $e$ to the ordinary integer zero, so the trace algebra has square-zero ideal spanned by $a,b$, as stated.

Suppose a coefficient refinement existed and apply its dimension map $\eps\mapsto1$. Write its dimension multiplication as
\[
 a^2=Aa,\qquad ab=Ba+Cb,
\]
where $A,B,C\in\mathbb N[q,q^{-1}]$. The absence of other outputs in $a^2$ is forced: an effective coefficient maps to the unsupported zero only when it is zero. Conversely the present coefficients at the $a$ and $b$ outputs of $ab$ force $B\ne0$ and $C\ne0$. At the $a$ coordinate, associativity gives
\[
 AB=AB+CB.
\]
Evaluating at $q=1$ contradicts $B(1)C(1)>0$. Thus there is no refinement. In fact the prescribed supports force these dimension coefficients to be positive constants, but the contradiction does not require that stronger observation.
\end{proof}
